\chapter*{Заключение}
\addcontentsline{toc}{chapter}{Заключение}
\label{ch:conclusion}


В рамках выпускной квалификационной работы была разработана интеллектуальная система HalalAI, объединяющая мобильное приложение, backend и LLM-сервис для генерации ответов, семантического поиска и работы с релевантными исламскими источниками на основе технологий RAG.


В процессе разработки был реализован backend на Spring Boot, обеспечивающий REST API, аутентификацию пользователей, работу с JWT-токенами и взаимодействие с LLM-сервисом. Также было разработано iOS-приложение на SwiftUI, включающее систему авторизации, AI-чат, экран чтения Корана, модуль анализа состава продуктов и систему уведомлений о времени намаза.


Особое внимание было уделено разработке и исследованию retrieval-подсистемы. В рамках экспериментального исследования были проверены гипотезы о влиянии использования RAG и дообучения embedding-моделей на качество ответов системы. Проведенные эксперименты показали, что использование retrieval-механизма действительно повышает обоснованность и релевантность ответов, а применение дообученных embedding-моделей улучшает качество семантического поиска.


Для обеспечения надежности проекта были написаны тесты для всех ключевых компонентов системы. Backend покрыт Unit и интеграционными тестами с использованием JUnit, Mockito и MockMvc. LLM-сервис протестирован с помощью pytest, было достигнуто полное покрытие строк тестируемого кода. Клиентское приложение покрыто Unit и UI-тестами, проверяющими основные пользовательские сценарии и бизнес-логику системы.


В ходе выполнения работы были решены следующие задачи:

\begin{enumerate}
	\item Проведен анализ предметной области и существующих решений для мусульманской аудитории.
	\item Сформулированы функциональные и нефункциональные требования к разрабатываемой системе.
	\item Разработана архитектура взаимодействия мобильного клиента, backend-сервиса и LLM-модуля.
	\item Спроектирована структура базы данных и retrieval-подсистема для семантического поиска.
	\item Разработан REST API для взаимодействия компонентов системы.
	\item Реализован LLM-сервис с использованием RAG-подхода.
	\item Разработано iOS-приложение, удовлетворяющее поставленным требованиям.
	\item Проведена интеграция компонентов системы и их тестирование.
\end{enumerate}

Таким образом, цель выпускной квалификационной работы была достигнута. Разработанная система продемонстрировала работоспособность выбранной архитектуры и эффективность применения retrieval-подсистемы и современных языковых моделей для решения задач, связанных с исламским информационным контекстом и анализом халяльности продуктов.
